\documentclass[10pt,a4paper]{amsart}
\usepackage[margin=19mm]{geometry}
\usepackage{amsmath,amssymb,mathtools}
\usepackage[scr=rsfs,cal=euler]{mathalfa}
\usepackage{xfrac}
\usepackage[unicode,hidelinks,bookmarks=false]{hyperref}
\newcommand{\ie}{i.e.\ }
\newcommand{\cf}{cf.\ }
\newcommand{\resp}{respectively\ }
\newcommand{\Subsection}{\subsection}
\providecommand{\nobreakdash}{\nobreak\hbox{-}}
\setlength{\emergencystretch}{2em}
\multlinegap0pt
\numberwithin{equation}{section}
\newtheorem{theo}{Theorem}[section]
\newtheorem{theorem}[theo]{Theorem}
\newtheorem{lemma}[theo]{Lemma}
\newtheorem{definition}[theo]{Definition}
\newtheorem{proposition}[theo]{Proposition}
\newtheorem{remark}[theo]{Remark}
\newtheorem{corollary}[theo]{Corollary}
\newtheorem{example}[theo]{Example}
\allowdisplaybreaks
\title[A note on sufficient conditions for the real Jacobian conjecture in the plane]{A quasi-homogeneous Newton-diagram sufficient condition for the real Jacobian conjecture in the plane (Theorem 1.7)}
\author{Xiuli Cen, Kangnong Hu, Yuzhou Tian}
\date{Source: Electronic Journal of Differential Equations 2026, No. 12, pp. 1--15; DOI 10.58997/ejde.2026.12; publisher version}
\begin{document}
\maketitle
\noindent\emph{Electronic Journal of Differential Equations}, Vol. 2026 (2026), No. 12, pp. 1--15.\newline
ISSN: 1072-6691. DOI \href{https://doi.org/10.58997/ejde.2026.12}{10.58997/ejde.2026.12}. This work is licensed under a CC BY 4.0 license.\par\smallskip
\noindent\textit{Editorial scope.} Counted claim: original Theorem 1.7 (the main theorem of the article, printed as Theorem 1.7 from the source label th1.8; source lines 164-177) together with its complete original proof, which the published article prints as the whole of Section 3 under the title Proof of Theorem 1.8 (source lines 379-695, containing the closing proof environment), delivered as the single primary proof body. Necessary complete context is retained uncounted: the whole of Section 1 up to the counted statement (source lines 66-163) with the real Jacobian conjecture and Theorems 1.1 to 1.6, the rest of Section 1 (source lines 179-346) with the figure that summarises the relationships between the sufficient conditions and the organisation paragraph, and the whole of Section 2 (source lines 200-346) with the Bendixson compactification, the Newton diagram and Lemmas 2.1 to 2.3, followed by the opening of Section 3 with Lemma 3.1 and its complete proof (source lines 347-377). Section 4 on the relationships between the existing theorems, with its examples, is not selected and is not used to inflate H. One printed cross-reference whose target is that undelivered Section 4 is delivered as a hyperlink whose visible text is the published section number 4. Because no content line of Sections 1 to 3 is dropped, every printed equation, lemma and theorem number of the delivered text coincides with the published PDF, in particular Lemma 2.1 to Lemma 2.3, Lemma 3.1 and Theorem 1.7. Figures: the retained spans contain five illustration asset lines, namely the single-panel figure of Section 1 and the two two-panel figures inside the counted proof; under the recorded bounded\_source\_comparison fidelity receipt only those asset lines are omitted, while every figure environment, caption, panel marker and cross-reference is retained, so the printed figure numbers used by the surrounding text are unchanged and no artwork is redistributed. Slips kept literal and recorded without mathematical repair: the source numbers its theorems with the label names th1.1, th1.2, th1.3, th1.5, th1.6, th1.7 and th1.8 while the shared counter prints them as Theorem 1.1 to Theorem 1.7, because Corollary 1.4 also consumes a number; the counted main theorem therefore prints as Theorem 1.7 although its label is th1.8, and the delivered text reproduces that exactly. No mathematics is written, repaired or re-derived.
\section{Introduction and main theorem}

In 1939, Keller \cite{bib 18} proposed the \emph{Jacobian conjecture} in
$\mathbb{C}^2$: If $F=(f,g):\mathbb{C}^2 \to \mathbb{C}^2$ is a polynomial map such 
that $ \det DF(x,y)$ is a non-zero constant for all $(x,y)\in \mathbb{C}^2$, 
then $ F $ is globally injective. See \cite{bib 4, bib 24, bib 26, bib 27} for 
further progress on this conjecture.

In 1983, Randall \cite{bib 22} raised the \emph{real Jacobian conjecture} in 
$\mathbb{R}^2$: If $F=(f,g):\mathbb{R}^2 \to \mathbb{R}^2$ is a polynomial map such 
that $ \det DF(x,y)\neq 0 $ for all $(x,y)\in \mathbb{R}^2$, then $ F $ is globally 
injective. As noted in \cite{bib 12,bib 26},  it is  beneficial to study the real 
Jacobian conjecture for understanding the Jacobian conjecture. Regrettably, 
Pinchuk \cite{bib 21} gave a counterexample to the real Jacobian conjecture in 
$\mathbb{R}^2$. As a result, a natural problem is to find sufficient conditions such 
that the real Jacobian conjecture holds. Based on algebraic and analytical methods, 
many results have been obtained, see \cite{bib 28, bib 6, bib 9, bib 13, bib 15, bib 17}. 

Cima et al.\ \cite{bib 11, bib 12} established a sufficient condition for the 
$n$-dimensional real Jacobian conjecture based on the structure of polynomial maps. 
When their finding is applied to the two-dimensional case, the result is as follows.

\begin{theorem}[\cite{bib 12}] \label{th1.1}
Assume that the polynomial map $F=(f,g):\mathbb{R}^2 \to \mathbb{R}^2$ satisfies 
$ F(0,0)=(0,0) $ and $ \det DF(x,y)\neq 0 $ for all $ (x,y)\in \mathbb{R}^2 $. 
If there exists a weight exponent $\mathbf{s}=(s_1,s_2)\in \mathbb{N}_{+}^{2}$
such that the higher $\mathbf{s}$-quasi-homogeneous part of $F$ only vanishes at 
the origin in $\mathbb{R}^2$, then $F$ is injective.
\end{theorem}

Sabatini \cite{bib 23} and Gavrilov \cite{bib 16} presented a global dynamical condition 
related to the real Jacobian conjecture in \(\mathbb{R}^2\). 
Based on their result, a variety of sufficient conditions for this conjecture 
have been obtained via the qualitative theory of dynamical systems, 
see \cite{bib 7, bib 8, bib 10, bib 20, bib 25}. 
One nice result is described as follows.

\begin{theorem}[\cite{bib 7}] \label{th1.2}
Assume that the polynomial map $F=(f,g):\mathbb{R}^2 \to \mathbb{R}^2$ satisfies $ F(0,0)=(0,0) $ and $ \det DF(x,y)\neq 0 $ for all $ (x, y)\in \mathbb{R}^2 $. If the higher homogeneous terms of the polynomials $ ff_{x}+gg_{x} $ and $ ff_{y}+gg_{y} $ do not have real linear factors in common, then $F$ is injective.
\end{theorem}

Braun et al.\ \cite{bib 7} also put forward an open problem that whether Theorem \ref{th1.1} implies Theorem \ref{th1.2}. However, they only solved this problem for some special cases. Liu and Tian \cite{bib 19} gave a positive answer and obtained a more general result for quasi-homogeneous type in $\mathbb{R}^n$. When their result is restricted to the two-dimensional case, the theorem is stated as follows.

\begin{theorem}[\cite{bib 19}] \label{th1.3}
Assume that the polynomial map $F=(f,g):\mathbb{R}^2 \to \mathbb{R}^2$ satisfies 
$ F(0,0)=(0,0) $ and $ \det DF(x,y)\neq 0 $ for all $ (x,y)\in \mathbb{R}^2 $.
 Let $H(x,y)=(f^2+g^2)/2$. If there exists a weight exponent 
$\mathbf{s}=(s_1,s_2)\in \mathbb{N}_{+}^{2}$ such that the higher 
$\mathbf{s}$-quasi-homogeneous part of vector field $(-\partial_{x}H,-\partial_{y}H)$ 
only vanishes at the origin in $\mathbb{R}^2$, then there exists a weight exponent 
$\Tilde{\mathbf{s}}=(\Tilde{s_1},\Tilde{s_2})\in \mathbb{N}_{+}^{2}$ such that the higher 
$\Tilde{\mathbf{s}}$-quasi-homogeneous part of $F$ only vanishes at the origin in 
$\mathbb{R}^2$.
\end{theorem}

Following Theorems \ref{th1.1} and \ref{th1.3}, a direct corollary can be obtained.

\begin{corollary}\label{th1.5}
Assume that the polynomial map $F=(f,g):\mathbb{R}^2 \to \mathbb{R}^2$ satisfies 
$ F(0,0)=(0,0) $ and $ \det DF(x,y)\neq 0 $ for all $ (x, y)\in \mathbb{R}^2 $. 
If there exists a weight exponent $\mathbf{s}=(s_1,s_2)\in \mathbb{N}_{+}^{2}$ such 
that the higher $\mathbf{s}$-quasi-homogeneous terms of the polynomials $ ff_{x}+gg_{x} $ 
and $ ff_{y}+gg_{y} $ only vanish at the origin, then $F$ is injective.
\end{corollary}

Tian and Cen \cite{bib 25} presented a  sufficient condition such that the 
two-dimensional real Jacobian conjecture holds through the Newton diagram. 
For the definition of the Newton diagram, see Section \ref{sec2}.

\begin{theorem}[\cite{bib 25}] \label{th1.6}
Consider a polynomial map $ F=(f, g):\mathbb{R}^2 \to \mathbb{R}^2 $ such that 
$ F(0,0)=(0,0) $ and $ \det DF(x,y)\neq 0 $ for all $ (x,y)\in \mathbb{R}^2 $. 
Denote by $ b(\mathcal{X}) $ the Bendixson compactification of the Hamiltonian vector 
field $ \mathcal{X}=(-ff_{y}-gg_{y}, ff_{x}+gg_{x})$, and by $ N(b(\mathcal{X})) $ 
the Newton diagram of $ b(\mathcal{X}) $. For each bounded edge of type 
$ \mathbf{t} $ in $ N(b(\mathcal{X})) $, if its associated Hamiltonian does not have 
any factor of the form $ v^{t_1}-\lambda u^{t_2} $ with 
$ \mathbf{t}=(t_{1},t_{2})\in N_{+}^{2} $ and 
$ \lambda \in \mathbb{R}\backslash \left\{0\right\} $, then $ F $ is injective.
\end{theorem}

It was also  proved that Theorem \ref{th1.6}  implies Theorem \ref{th1.2}.

\begin{theorem}[\cite{bib 25}] \label{th1.7}
Consider a polynomial map $ F=(f, g):\mathbb{R}^2 \to \mathbb{R}^2 $ such that 
$ F(0,0)=(0,0) $ and $ \det DF(x,y)\neq 0 $ for all $ (x,y)\in \mathbb{R}^2 $. 
Denote by $ b(\mathcal{X}) $ the Bendixson compactification of the Hamiltonian vector 
field $ \mathcal{X}=(-ff_{y}-gg_{y}, ff_{x}+gg_{x})$, and by $ N(b(\mathcal{X})) $ 
the Newton diagram of $ b(\mathcal{X}) $. If the higher homogeneous terms of the 
polynomials $ ff_{x}+gg_{x} $ and $ ff_{y}+gg_{y} $ do not have real linear factors 
in common, then for each bounded edge of type $ \mathbf{t} $ in $ N(b(\mathcal{X})) $, 
its associated Hamiltonian does not have any factor of the form 
$ v^{t_1}-\lambda u^{t_2} $ with $ \mathbf{t}=(t_{1},t_{2})\in N_{+}^{2} $ and 
$ \lambda \in \mathbb{R}\backslash \left\{0\right\}$.
\end{theorem}

In this article, we generalize Theorem \ref{th1.7} to the quasi-homogeneous case.


\begin{theorem}\label{th1.8}
Consider a polynomial map $ F=(f,g):\mathbb{R}^2 \to \mathbb{R}^2 $ such that 
$ F(0,0)=(0,0) $ and $ \det DF(x,y)\neq 0 $ for all $ (x,y)\in \mathbb{R}^2 $. 
Denote by $ b(\mathcal{X}) $ the Bendixson compactification of the Hamiltonian vector 
field $ \mathcal{X}=(-ff_{y}-gg_{y}, ff_{x}+gg_{x}) $, and by $ N(b(\mathcal{X})) $ 
the Newton diagram of $ b(\mathcal{X}) $. If there exists a weight exponent 
$\mathbf{s}=(s_1,s_2)\in \mathbb{N}_{+}^{2}$ such that the higher 
$\mathbf{s}$-quasi-homogeneous terms of the polynomials $ ff_{x}+gg_{x} $ and 
$ ff_{y}+gg_{y} $ only vanish at the origin, then for each bounded edge of type 
$ \mathbf{t} $ in $ N(b(\mathcal{X})) $, its associated Hamiltonian does not have any 
factor of the form $ v^{t_1}-\lambda u^{t_2} $ with 
$ \mathbf{t}=(t_{1},t_{2})\in N_{+}^{2} $ and 
$ \lambda \in \mathbb{R}\backslash \left\{0\right\} $.
\end{theorem}

Theorem \ref{th1.8} tells us that the condition in Theorem \ref{th1.6} is weaker than 
the condition in Corollary \ref{th1.5}. Now, we are intrigued by the relationships 
between the sufficient conditions for Theorems \ref{th1.1}, \ref{th1.2}, \ref{th1.6} 
and Corollary \ref{th1.5}.  Combining with the existing results, we illustrate the 
relationships by some examples, as shown in Figure \ref{Fig.main2}.

\begin{figure}[htb]
\centering
\caption{Relationships between the sufficient conditions for Theorems \ref{th1.1}, 
\ref{th1.2}, \ref{th1.6}, and Corollary \ref{th1.5}.}
\label{Fig.main2} 
\end{figure}

This article is organized as follows. 
In Section \ref{sec2}, we provide some preliminary results. 
In Section \ref{sec3}, we give a proof of Theorem \ref{th1.8}. 
Finally, some examples are presented to show the relationships between the sufficient 
conditions for Theorems \ref{th1.1}, \ref{th1.2}, \ref{th1.6} and 
Corollary \ref{th1.5} in Section \href{https://ejde.math.txstate.edu/Volumes/2026/12/cen.pdf}{4}.

\section{Preliminaries}\label{sec2}

\subsection{Bendixson compactification}
Consider a Hamiltonian vector field
\begin{equation}\label{eq1-1}
\mathcal{X}=(-ff_{y}-gg_{y}, ff_{x}+gg_{x}) := (\Phi(x,y), \Psi(x,y)).
\end{equation}
Let $ d_0 = \max \{\deg \Phi(x,y), \deg \Psi(x,y)\} $.  
We introduce the Bendixson transformation, 
\[ %\label{eq2-1}
 x =\frac{u}{u^{2}+v^{2}},\quad y =\frac{v}{u^{2}+v^{2}},
\]
and denote by $ b(\mathcal{X}) $ the Bendixson compactification of the Hamiltonian 
vector field $\mathcal{X}$ given in \eqref{eq1-1}. Then the explicit expression for
$ b(\mathcal{X}) $ is
\begin{equation}\label{eq2-2}
\begin{gathered}
\dot u  = (u^{2}+v^{2})^{d_0}
\Big[(v^{2}-u^{2})\Phi\Big(\frac{u}{u^{2}+v^{2}}, \frac{v}{u^{2}+v^{2}}\Big)
-2uv\Psi \Big(\frac{u}{u^{2}+v^{2}}, \frac{v}{u^{2}+v^{2}}\Big) \Big], \\
\dot v  = (u^{2}+v^{2})^{d_0}
\Big[(u^{2}-v^{2})\Psi \Big(\frac{u}{u^{2}+v^{2}}, \frac{v}{u^{2}+v^{2}}\Big)
-2uv\Phi\Big(\frac{u}{u^{2}+v^{2}}, \frac{v}{u^{2}+v^{2}}\Big) \Big].
\end{gathered}
\end{equation}
For more details about Bendixson compactification, see \cite{bib 3, bib 14}.

Let $f(x,y)=\sum_{i=1}^{n}f_{i}(x,y)$, $g(x,y)=\sum_{j=1}^{m}g_{j}(x,y)$ and 
$d=\max\{n,m\}$, where $f_{i}(x,y)$ and $g_{j}(x,y)$ are homogeneous polynomials of 
degree $i$ and $j$, respectively. From equation \eqref{eq2-2}, the expression for 
$b(\mathcal{X})$ can be written as
\begin{gather*} %\label{eq2-3}
\begin{aligned}
\dot u &=  \sum_{i=1}^{d}\sum_{j=1}^{d}(u^{2}+v^{2})^{2d-i-j}[(u^{2}-v^{2})(f_{i}(u,v)f_{jy}(u,v)+g_{i}(u,v)g_{jy}(u,v)) \\
&\quad -2uv(f_{i}(u,v)f_{jx}(u,v)+g_{i}(u,v)g_{jx}(u,v))],
\end{aligned} \\
\begin{aligned}
\dot v  &=  \sum_{i=1}^{d}\sum_{j=1}^{d}(u^{2}+v^{2})^{2d-i-j}[(u^{2}-v^{2})(f_{i}(u,v)f_{jx}(u,v)+g_{i}(u,v)g_{jx}(u,v)) \\
&\quad +2uv(f_{i}(u,v)f_{jy}(u,v)+g_{i}(u,v)g_{jy}(u,v))].
\end{aligned}
\end{gather*}
Note that $f_{jx}(u,v)=f_{ju}(u,v),\ f_{jy}(u,v)=f_{jv}(u,v),
 g_{jx}(u,v)=g_{ju}(u,v)$ and $g_{jy}(u,v)=g_{jv}(u,v)$. 
The vector field $b(\mathcal{X})$ can be rewritten as the sum of its homogeneous 
components
\begin{equation}\label{eq2-4}
b(\mathcal{X})=\frac{1}{2}\sum_{i=1}^{d}\mathcal{F}_{4d+1-2i}
+\sum_{1\leq i<j\leq d}\mathcal{F}_{4d+1-i-j},
\end{equation}
where
\begin{equation}\label{eq2-5}
 \mathcal{F}_{4d+1-i-j}
= \begin{pmatrix}
 (u^{2}+v^{2})^{2d-i-j}[(u^{2}-v^{2})\partial_{v}(f_{i}f_{j}+g_{i}g_{j})
  -2uv\partial_{u}(f_{i}f_{j}+g_{i}g_{j})] \\
 (u^{2}+v^{2})^{2d-i-j}[(u^{2}-v^{2})\partial_{u}(f_{i}f_{j}+g_{i}g_{j})
  +2uv\partial_{v}(f_{i}f_{j}+g_{i}g_{j})]
\end{pmatrix}^T.
\end{equation}

\subsection{Newton diagram}

Let $\mathbf{s} = (s_{1},s_{2})\neq \boldsymbol{0} $ with $s_{1}, s_{2}\in\mathbb{N}$ 
coprime. We say that a polynomial $R(x,y)$ is quasi-homogeneous of weighted degree 
$k$ with respect to weight exponent $\mathbf{s} = (s_{1},s_{2})$ if 
$R(\lambda^{s_{1}}x,\lambda^{s_{2}}y)=\lambda^{k}R(x,y)$ for all $\lambda >0$. 
Denote by $\mathcal{P}_{k}^{\mathbf{s}}$ the vector space of quasi-homogeneous 
polynomials of weighted degree $k$ with respect to weight exponent 
$\mathbf{s} = (s_{1},s_{2})$. If $P_{k+s_{1}}\in \mathcal{P}_{k+s_{1}}^{\mathbf{s}}$ and 
$Q_{k+s_{2}}\in \mathcal{P}_{k+s_{2}}^{\mathbf{s}}$, then the vector field 
$\mathscr{X}_{k}=(P_{k+s_{1}},Q_{k+s_{2}})$ is called quasi-homogeneous of weighted 
degree $k$ with respect to weight exponent $\mathbf{s} = (s_{1},s_{2})$. 
Denote by $\mathcal{Q}_{k}^{\mathbf{s}}$ the vector space of the quasi-homogeneous 
polynomial vector fields of weighted degree $k$ with respect to weight exponent 
$\mathbf{s} = (s_{1},s_{2})$.

 We introduce the Newton diagram briefly, see \cite{bib 2, bib 5} for more details.
Consider the vector field $\mathscr{X}= (P(x,y), Q(x,y))$. 
Let $yP(x,y)=\sum a_{ij}x^{i}y^{j}$, $xQ(x,y)=\sum b_{ij}x^{i}y^{j}$ and 
$R(x,y)=\sum c_{ij}x^{i}y^{j}$. The \emph{support} of $\mathscr{X}$ is 
\[
\operatorname{supp}(\mathscr{X})=\{(i,j)\mid(a_{ij},b_{ij})\neq (0,0)\}
\subseteq \mathbb{R}^2.
\]
The vector $(a_{ij},b_{ij})$ is called the \emph{vector coefficient} corresponding to 
$(i,j)$ in the support. Similarly, the set
\[
\operatorname{supp}(R)=\{(i,j)\mid c_{ij}\neq 0\}\subseteq \mathbb{R}^2
\]
is called the \emph{support} of the polynomial $R(x,y)$.

The boundary of the convex hull for the set
$$
\cup_{(i,j)\in \operatorname{supp}(\mathscr{X})}((i,j)+\mathbb{R}_{+}^{2})
$$
includes two open rays and a polygonal line that may be just a single point. 
This polygonal line with the rays that do not lie on a coordinate axis, when they exist, 
is called the \emph{Newton diagram} of the vector field $\mathscr{X}$. 
The Newton diagram of the polynomial $R(x,y)$ can be defined similarly.

The line segments and their endpoints of the polygonal line are the \emph{edges} 
and \emph{vertices} of the Newton diagram, respectively. 
A vertex is classified as an \emph{exterior vertex} if it lies on a coordinate axis; 
otherwise, it is called an \emph{inner vertex}. The \emph{exponent} of a 
\emph{bounded edge} $l$ in the Newton diagram is defined as a positive rational number 
$t_2/t_1$, which equals the reciprocal of the absolute value of the edge's slope. 
The pair $t=(t_1,t_2)$ is called the \emph{type} of the edge $l$.

For each bounded edge of type $\mathbf{t}=(t_1,t_2)$, we can define a quasi-homogeneous 
Hamiltonian vector field 
$\mathscr{X}_{k}=(P_{k+t_{1}},Q_{k+t_{2}})\in \mathcal{Q}_{k}^{\mathbf{t}}$. 
The associated Hamiltonian for this vector field is as follows
\begin{equation}\label{eq2-6}
h_{k+|\mathbf{t}|}=\frac{1}{k+|\mathbf{t}|}(t_{1}xQ_{k+t_2}-t_{2}yP_{k+t_1)}
\in \mathcal{P}_{k+|\mathbf{t}|}^{\mathbf{t}}, \quad |\mathbf{t}|=t_1+t_2,
\end{equation}
see \cite{bib 1} for more details.

\subsection{Properties of the Newton diagram}

\begin{lemma}[\cite{bib 25}] \label{lem2.1}
Assume that $R(x,y)\in \mathcal{P}_{k}^{\textbf{t}}$ and the vector 
field  $\mathscr{X}_{k}=(P_{k+t_{1}},Q_{k+t_{2}})\in \mathcal{Q}_{k}^{\textbf{t}}$. 
Then $\operatorname{supp}(\mathscr{X}_{k})$ lies on the straight line 
$t_1x+t_2y=k+|\textbf{t}|$, and $\operatorname{supp}(R)$ lies on the straight line
$t_1x+t_2y=k$.
\end{lemma}

\begin{lemma}[\cite{bib 25}] \label{lem2.2}
Consider the vector field $\mathscr{X}=\sum \mathscr{X}_{i}$ with 
$\mathscr{X}_{i}\in \mathcal{Q}_{i}^{\textbf{t}}$. Let $V_{1i}$ and $V_{2i}$ be the 
vertices of the Newton diagram of $\mathscr{X}_{i}$. Then the vertices of the Newton 
diagram of $\mathscr{X}$ are contained in
the set $\cup_{i}\{V_{1i},V_{2i}\}\subset \operatorname{supp}(\mathscr{X})$.
\end{lemma}

\begin{lemma}[\cite{bib 25}] \label{lem2.3}
 Consider the vector field
\[
\mathcal{Y}=((x^2-y^2)\partial_{y}\mathcal{H}-2xy\partial_{x}\mathcal{H}, 
(x^2-y^2)\partial_{x}\mathcal{H}+2xy\partial_{y}\mathcal{H}),
\]
where $\mathcal{H}=\mathcal{H}(x,y)$ is a homogeneous polynomial of degree $k$ in 
the variables $x$ and $y$. Then the Newton diagram of the vector field $\mathcal{Y}$ 
and the polynomial $(x^2+y^2)\mathcal{H}(x,y)$ have the same vertices.
\end{lemma}


\section{Proof of Theorem \ref{th1.8}}\label{sec3}

We first give the following lemma.

\begin{lemma}\label{lem3.2}
Assume that the higher $\mathbf{s}$-quasi-homogeneous term of the polynomial $H(x,y)$ 
with respect to the weight exponent $\mathbf{s}$ = $(s_1,s_2)$ has the  form
\[
H^{\mathbf{s}}(x,y)=a_1x^{k_1}+a_2y^{k_2}+\sum_{i\geq 1, j\geq 1} c_{ij}x^{i}y^{j},
\]
where $a_1>0$, $a_2>0$, $c_{ij}\in \mathbb{R}, $ $k_1\in \mathbb{N}_+$, 
$k_2\in \mathbb{N}_+$, and the equality $k_1s_1=k_2s_2=is_1+js_2$ holds. 
If $s_1<s_2$, then the higher homogeneous term of $H(x,y)$ is $a_1x^{k_1}$, while 
the highest degree term with respect to $y$ in $H(x,y)$ is $a_2y^{k_2}$, and the 
inequality $k_2<i+j<k_1$ holds.
\end{lemma}

\begin{proof}
Since $k_1s_1=k_2s_2$ and $s_1<s_2$, then $k_1>k_2$. 
Let $a_{m-l,l}x^{m-l}y^{l}\ (l\geq 1)$ be a term of degree $m$ in $H(x,y)$. 
If $m\geq k_1$, we have $(m-l)s_1+ls_2> ms_1\geq k_1s_1$. This contradicts the fact 
that $a_1x^{k_1}$ is the higher $\mathbf{s}$-quasi-homogeneous term of $H(x,y)$. 
Therefore, $m<k_1$ and the higher homogeneous term of $H(x,y)$ can only be $a_1x^{k_1}$. 
Similarly, we can verify that the highest degree term with respect to $y$ in $H(x,y)$ 
is $a_2y^{k_2}$.

Now, we prove $i+j > k_2$ by contradiction. If $i+j\leq k_2$, then we have 
$is_1+js_2<(i+j)s_2\leq k_2s_2$. This is in contradiction with $is_1+js_2=k_2s_2$. 
Hence, $k_2<i+j<k_1$.
This completes the proof.
\end{proof}

\begin{proof}[Proof of Theorem \ref{th1.8}]
Let $f(x,y)$ and $g(x,y)$  be the polynomials of degree $n$ and $m$, respectively. 
Assume that $f(x,y)=\sum_{i=1}^{n_{\mathbf{s}}}f^{\mathbf{s}}_{i}(x,y)$, 
$g(x,y)=\sum_{j=1}^{m_{\mathbf{s}}}g^{\mathbf{s}}_{j}(x,y)$, $d=\max\{n,m\}$ and 
$d_{\mathbf{s}}=\max\{n_{\mathbf{s}},m_{\mathbf{s}}\}$, where 
$f^{\mathbf{s}}_{i}(x,y)$ and $g^{\mathbf{s}}_{j}(x,y)$ are quasi-homogeneous 
polynomials of weighted degree $i$ and $j$ with respect to weight exponent 
$\mathbf{s}=(s_1,s_2)$, respectively. Then
\[
H(x,y)=\frac{f^2(x,y)+g^2(x,y)}{2}
=\frac{1}{2}\sum_{j=2}^{2d_{\mathbf{s}}}H^{\mathbf{s}}_{j}(x,y),
\]
where $H^{\mathbf{s}}_j(x,y)$ is a quasi-homogeneous polynomial of weighted degree 
$j$ with respect to weight exponent $\mathbf{s}$ = $(s_1,s_2)$.

According to the condition of Theorem \ref{th1.8}, there exist the 
$\mathbf{s}$-quasi-homogeneous terms  $H^{\mathbf{s}}_{r_1}(x,y)$ and 
$H^{\mathbf{s}}_{r_2}(x,y)$ in  $H(x,y)$ such that
$\partial_{x}H^{\mathbf{s}}_{r_1}$ and $\partial_{y}H^{\mathbf{s}}_{r_2}$ only 
vanish at the origin. Without loss of generality, we can assume that 
$2d_{\mathbf{s}}=r_1\geq r_2$ and $s_1<s_2$. Consequently, our proof can be divided into 
the following two cases.
\smallskip

\noindent\textbf{Case 1:} $r_1=r_2=2d_{\mathbf{s}}$.
In this case, the higher $\mathbf{s}$-quasi-homogeneous terms of the polynomials 
$ ff_{x}+gg_{x} $ and $ ff_{y}+gg_{y} $ only vanish at $ (0,0) $ if and only if 
$\partial_{x}H^{\mathbf{s}}_{2d_{\mathbf{s}}}$ and 
$\partial_{y}H^{\mathbf{s}}_{2d_{\mathbf{s}}}$ only vanish at $ (0,0)$. 
This means that $(0,0)$ is the unique real zero of 
$H^{\mathbf{s}}_{2d_{\mathbf{s}}}(x,y)$, see \cite[Corollary 1]{bib 19}.

Suppose that $H^{\mathbf{s}}_{2d_{\mathbf{s}}}(x,y)$ has the general form
\begin{equation}\label{eq3-1}
H^{\mathbf{s}}_{2d_{\mathbf{s}}}(x,y)
=Ax^{k_1}+By^{k_2}+\sum_{i_0s_1+j_0s_2=2d_{\mathbf{s}}, 
i_0, j_0\geq1}C_{i_0j_0}x^{i_0}y^{j_0},
\end{equation}
where
\begin{equation}\label{eq3-2}
k_1s_1=k_2s_2=i_0s_1+j_0s_2=2d_{\mathbf{s}}.
\end{equation}
Since $H^{\mathbf{s}}_{2d_{\mathbf{s}}}(x,y)
=\left((f^{\mathbf{s}}_{d_{\mathbf{s}}})^2+(g^{\mathbf{s}}_{d_{\mathbf{s}}})^2
\right)/2\geq 0$, 
we conclude that $H^{\mathbf{s}}_{2d_{\mathbf{s}}}(0,y)=By^{k_2}\geq 0$ and 
$H^{\mathbf{s}}_{2d_{\mathbf{s}}}(x,0)=Ax^{k_1}\geq 0$. If either $A=0$ or $B=0$, 
then it would contradict the fact that $(0,0)$ is the unique real zero of 
$H^{\mathbf{s}}_{2d_{\mathbf{s}}}(x,y)$. Therefore, $A>0$, $B>0$ and $k_1$ and 
$k_2$ are even integers.
According to Lemma \ref{lem3.2}, the higher homogeneous term in $H(x,y)$ is $Ax^{k_1}$. 
Thus, $k_1=2d$, and $H_{2d}(x,y)=Ax^{2d}$. Moreover, let $k_2=2d_2$,
it follows from the equality \eqref{eq3-2} that $s_1=d_{\mathbf{s}}/d$, 
$s_2=d_{\mathbf{s}}/d_2$, and
\begin{equation}\label{eq3-3}
i_0d_2+j_0d=2dd_{2}.
\end{equation}

It is clear that the Hamiltonian $H(x,y)$ has the  form
\[ %\label{eq3-4}
 H(x,y)= \frac{1}{2}\left(H_{2d}(x,y)+\cdots+H_{k}(x,y)+\cdots
 +H_{2d_2}(x,y)+\cdots+H_{2}(x,y)\right)
= \frac{1}{2}\sum_{l=2}^{2d}H_l(x,y),
\]
where $2d_2<k<2d$, and $H_l(x,y)$ is a homogeneous polynomial of degree $l$ for
$l=2,3,\dots,2d$. By Lemma \ref{lem3.2}, the degree of each term with respect to $y$ 
in $H_k(x,y)$ is less than $2d_2$, and the polynomial $H_{2 d_2}(x,y)$ takes the form
\begin{equation}\label{Hd2}
H_{2d_2}(x,y)=By^{2d_2}+\sum_{i+j=2d_2, j<2d_2}b_{ij}x^{i}y^{j}.
\end{equation}
Moreover, the higher $\mathbf{s}$-quasi-homogeneous terms in 
$H^{\mathbf{s}}_{2d_{\mathbf{s}}}(x,y)$ cannot appear in all $H_l(x,y)$ with 
$2\leq l<2d_2$.

By equation \eqref{eq2-2}, the expression of $b(\mathcal{X})$ is given by
\begin{equation}\label{eq3-5}
\begin{gathered}
\dot u  = \frac{1}{2}\sum_{l=2}^{2d}(u^{2}+v^{2})^{2d-l}[(u^{2}-v^{2})H_{l,v}(u,v)-2uvH_{l,u}(u,v)], \\
\dot v  = \frac{1}{2}\sum_{l=2}^{2d}(u^{2}+v^{2})^{2d-l}[(u^{2}-v^{2})H_{l,u}(u,v)+2uvH_{l,v}(u,v)],
\end{gathered}
\end{equation}
where $H_{l,u}(u,v)$ and $H_{l,v}(u,v)$ represent the partial derivatives of 
$H_l(u,v)$ with respect to $u$ and $v$, respectively. Let
\begin{equation}\label{eq3-6}
\mathcal{F}_{4d+1-l}
=\begin{pmatrix}
 (u^{2}+v^{2})^{2d-l}[(u^{2}-v^{2})H_{l,v}(u,v)-2uvH_{l,u}(u,v)] \\
 (u^{2}+v^{2})^{2d-l}[(u^{2}-v^{2})H_{l,u}(u,v)+2uvH_{l,v}(u,v)]
\end{pmatrix}^T\in \mathcal{Q}_{4d+1-l}^{(1,1)}
\end{equation}
for $l=2,3,\dots,2d$. Then
$$
\operatorname{supp}(b(\mathcal{X}))
=\cup_{l=2}^{2d}\operatorname{supp}(\mathcal{F}_{4d+1-l}).
$$

By Lemma \ref{lem2.1}, the support of $\mathcal{F}_{4d+1-l}$ lies on the straight 
line $\mathcal{L}_{4d+1-l}:x+y=4d+2-l$ for $l=2,3,\dots,2d$. Additionally, from 
Lemma \ref{lem2.3}, we know that  the Newton diagrams of $\mathcal{F}_{4d+1-l}$ 
and $(u^2+v^2)^{2d+1-l}H_{l}(u,v)$ have the same vertices for $l=2,3,\dots,2d$.

Recall that $H_{2d}(x,y)=Ax^{2d}$. Thus, for $l=2d$, we have 
$(u^2+v^2)H_{2d}(u,v)=Au^{2d}(u^2+v^2)$. Consequently, the vertices of the Newton 
diagram of $\mathcal{F}_{2d+1}$ are
\[
V_{1,2d}=(2(d+1),0)\quad \text{and} \quad V_{2,2d}=(2d,2).
\]
For $l=2d_2$, from the expression of $H_{2d_2}(x,y)$ given by \eqref{Hd2}, 
the Newton diagram of $\mathcal{F}_{4d+1-2d_2}$ has an exterior vertex
\[
V_{2,2d_2}=(0,2(2d-d_2+1)).
\]
Note that $V_{2,2d}$ and $V_{2,2d_2}$ lie on the straight line
\begin{equation}\label{l1}
\mathcal{L}_0:(2d-d_2)x+dy=2d(2d-d_2+1).
\end{equation}
For $2d_2<l<2d$, two cases are considered.

(i) If $H_{l}(x,y)$ includes the higher $\mathbf{s}$-quasi-homogeneous term 
$C_{i_{0}j_{0}}x^{i_0}y^{j_0}$,  then it can be expressed as
$$
H_{l}(x,y)= C_{i_{0}j_{0}}x^{i_0}y^{j_0}+\sum_{i+j=l, j<j_0}h_{ij}x^iy^j,
$$
where $i_0=l-j_0$. Consequently, one of the vertices of the Newton diagram of 
$\mathcal{F}_{4d+1-l}$ is given by
\[
V_{2,l}=(i_0, 2(2d+1-l)+j_0).
\]
Recall that $i_0d_2+j_0d=2d d_2$, see \eqref{eq3-3}. It is easy to verify that
\[ %\label{eq3-8}
(2d-d_2)i_0+d(2(2d+1-l)+j_0)=2d(2d+1)-(i_0d_2+j_0d)=2d(2d-d_2+1).
\]
This implies that the vertex $V_{2,l}$ lies on the line $\mathcal{L}_0$ given by 
\eqref{l1}.

(ii) If $H_{l}(x,y)$ does not contain the higher $\mathbf{s}$-quasi-homogeneous term 
$C_{i_{0}j_{0}}x^{i_0}y^{j_0}$, then
$$
H_{l}(x,y)= \sum_{i+j=l, j<j_0}h_{ij}x^iy^j.
$$
Let $j^*$ be the maximum degree of $H_l(x,y)$ with respect to $y$ and $i^*=l-j^*$. 
Therefore, $i^*d_2+j^*d<2d d_2$. As a result, one of the vertices of the Newton 
diagram of $\mathcal{F}_{4d+1-l}$ is
\[
V_{2,l}=(i^*, 2(2d+1-l)+j^*).
\]
Note that
\[
(2d-d_2)i^*+d(2(2d+1-l)+j^*)=2d(2d+1)-(i^*d_2+j^*d)>2d(2d-d_2+1).
\]
This shows that the vertex $V_{2,l}$ is located to the right of the line $\mathcal{L}_0$.

Owing to the above analysis, we obtain the configuration of the support for the vector 
field $b(\mathcal{X})$, see Figure \ref{Fig.main}(a).  By Lemma \ref{lem2.2}, 
the Newton diagram of $b(\mathcal{X})$ has two exterior vertices $V_{1,2d}$ and 
$V_{2,2d_2}$, and an inner vertex $V_{2,2d}$, and  it  includes two edges of type 
$(2 d-d_2, d)$ and $(1,1)$, as illustrated in Figure \ref{Fig.main}(b). 
The vector fields associated to the vertices $V_{1,2 d}, V_{2,2 d}$ and $V_{2,2 d_2}$ 
are $(0, Adu^{2d+1})$, $(-2Ad u^{2 d}v, -Ad u^{2 d-1}v^{2})$ and
$(-Bd_2 v^{4 d-2 d_2+1},0)$, respectively.

\begin{figure}[htp]
\centering
(a) \hfil (b)
\caption{ (a) Configuration of the support of $b(\mathcal{X})$,
(b)  Newton diagram for $b(\mathcal{X})$.}
\label{Fig.main}
\end{figure}

For the vector field $b(\mathcal{X})$, the lowest-degree homogeneous terms with weight 
exponent $(1,1)$ are
\[
\frac{1}{2}\mathcal{F}_{2d+1}
=\frac{1}{2}\begin{pmatrix}
& (u^{2}-v^{2})(H_{2d,v}(u,v)-2uvH_{2d,u}(u,v)) \\
& (u^{2}-v^{2})(H_{2d,u}(u,v)+2uvH_{2d,v}(u,v))
\end{pmatrix}^T
= \begin{pmatrix}
& -2Adu^{2d}v \\
& Ad(u^2-v^2)u^{2d-1}
\end{pmatrix}^T.
\]
By \eqref{eq2-6}, the Hamiltonian associated with the edge of type $(1,1)$ is
\[
h_{2(d+1)}^{(1,1)}=\frac{Adu^{2d}}{2(d+1)}(u^2+v^2).
\]
It is clear that $h_{2d+2}^{(1,1)}$ does not contain any factor of the form 
$v-\lambda u$ with $\lambda \in \mathbb{R}\backslash\{0\}$.

The lowest-degree quasi-homogeneous terms with weight exponent $(2d-d_2, d)$ for the 
vector field $b(\mathcal{X})$ are
\[
\begin{pmatrix}
 -2Adu^{2d}v-\frac{1}{2}\sum_{i_0d_2+j_0d=2d d_{2}}C_{i_0j_0}(j_0+2i_0)u^{i_0}
 v^{4d-2i_0-j_0+1}-Bd_2v^{4d-2d_2+1}  \\
 -Adu^{2d-1}v^{2}-\frac{1}{2}\sum_{i_0d_2+j_0d=2d d_{2}}i_0C_{i_0j_0}u^{i_0-1}
 v^{4d-2i_0-j_0+2}
\end{pmatrix}^T.
\]
By \eqref{eq2-6}, the Hamiltonian associated with the edge of type $(2d-d_2, d)$ is
\[
h_{2d(2d-d_2+1)}^{(2d-d_2,d)}
=\frac{d_2v^{2}}{2(2d-d_2+1)}
\Big(Au^{2d}+Bv^{4d-2d_2}+\sum_{i_0d_2+j_0d=2d d_{2}}C_{i_0j_0}u^{i_0}v^{4d-2i_0-j_0}\Big).
\]
Note that $\mathbf{s}=(s_1,s_2)=(d_2, d)$ and
$$
H^{\mathbf{s}}_{2d_{\mathbf{s}}}(x,y)=Ax^{2d}
+By^{2d_2}+\sum_{i_0d_2+j_0d=2d d_{2}, i_0, j_0\geq1}C_{i_0j_0}x^{i_0}y^{j_0}\quad
 \text{(see \eqref{eq3-1})}
 $$
only vanishes at the origin. Thus the Hamiltonian
\[
h_{2d(2d-d_2+1)}^{(2d-d_2,d)}
=\frac{d_2v^{2}}{2(2d-d_2+1)}H^{\mathbf{s}}_{2d_{\mathbf{s}}}
\left(u,v^{\frac{2d-d_2}{d_2}}\right)
\]
has no factor of the form $v^{2d-d_2}-\lambda u^{d}$ with $\lambda\neq 0$. 
Otherwise, the higher $\mathbf{s}$-quasi-homogeneous term 
$H^{\mathbf{s}}_{2d_{\mathbf{s}}}(x,y)$ will vanish at $y^{d_2}-\lambda x^d=0$.
\smallskip

\noindent\textbf{Case 2:}
$r_1=2d_{\mathbf{s}}>r_2$.
In this case, 
\begin{equation}\label{eq3-7}
H(x,y)=\frac{1}{2}\left(H^{\mathbf{s}}_{2d_{\mathbf{s}}}(x)+\dots
+H^{\mathbf{s}}_{k}(x)+\dots+H^{\mathbf{s}}_{r_2+1}(x)+H^{\mathbf{s}}_{r_2}(x,y)
+\dots+H^{\mathbf{s}}_{2}(x,y)\right),
\end{equation}
where $H^{\mathbf{s}}_k(x)=A_kx^{a_k}$, $k=r_2+1, \ldots, 2d_{\mathbf{s}}$, and 
$H^{\mathbf{s}}_{k}(x,y)$, $k=2,3,\ldots, r_2$ are quasi-homogeneous polynomials of 
weighted degree $k$ with weight exponent $\mathbf{s}=(s_1,s_2)$. 
The polynomial $H^{\mathbf{s}}_{r_2}(x,y)$ has the  form
\[
H^{\mathbf{s}}_{r_2}(x,y)=By^{b}+\sum_{is_1+js_2
=r_2}b_{ij}x^{i}y^{j}+A_{r_2}x^{a_{r_2}},
\]
where $b_{ij}\in \mathbb{R}$ and $A_{r_2}\in \mathbb{R}$.  It follows from
\[
H^{\mathbf{s}}_{2d_{\mathbf{s}}}(x)=
\frac{1}{2}\left(\left(f^{\mathbf{s}}_{d_{\mathbf{s}}}\right)^2
+\left(g^{\mathbf{s}}_{d_{\mathbf{s}}}\right)^2\right)
=\frac{1}{2}A_{2d_{\mathbf{s}}}x^{a_{2d_{\mathbf{s}}}}\geq 0
\]
that $A_{2d_{\mathbf{s}}}>0$ and $a_{2d_{\mathbf{s}}}=2d$. 
By Lemma \ref{lem3.2}, it is easy to verify that the term  $By^b$ is the highest 
degree term in $H(x,y)$ with respect to $y$. Therefore,
\[
\frac{f^{2}(0,y)+g^{2}(0,y)}{2}=H(0,y)=\frac{1}{2}(By^b
+\dots+H^{\mathbf{s}}_{2}(0,y))\geq 0.
\]
This implies that $B\geq 0$ and $b$ is even. Thus, $r_2=bs_2$ is also even. 
Since $H^{\mathbf{s}}_{2d_{\mathbf{s}}}(x)$ and $H^{\mathbf{s}}_{r_2}(x,y)$ only 
vanish at $(0,0)$,  we have $B>0$.  For convenience, we denote $a_{r_2}=d_1$ and 
$b=2d_2$.

By Lemma \ref{lem3.2},  $H(x,y)$  can be arranged in descending homogeneous order as 
follows
\begin{align*}
    H(x,y)&=\frac{1}{2}\left(H_{2d}(x)\!+\!\dots\!+\!H_{a_k}(x)\!+\!\dots\!+\!H_{d_1}(x)\!+\!\dots\!+\!H_{2d_2}(x,y)\!+\!\dots\!+\!H_2(x,y)\right)\\
&=\frac{1}{2}\left(A_{2d_{\mathbf{s}}}x^{2d}+\dots+A_{k}x^{a_{k}}+\dots+A_{r_2}x^{d_1}+\dots+By^{2d_2}+\dots+H_2(x,y)\right).
\end{align*}
Similar to the computation in Case 1, we have that  
$(u^2+v^2)^{2d+1-l}H_{l}(u,v)=(u^2+v^2)^{2d+1-l}A_{ls_1}u^l$, and the vertices of the 
Newton diagram of $\mathcal{F}_{4d+1-l}$ are
\[
V_{1,l}=(2(2d+1)-l,0) \quad \text{and} \quad V_{2,l}=(l,2(2d+1-l))
\]
for $l=d_1, d_1+1, \ldots, 2d$.
Note that $V_{2,l}$ lies on the straight line $L_1:2x+y=2(2d+1)$.

For $l=2d_2$, the Newton diagram of $\mathcal{F}_{4d+1-2d_2}$ has an exterior vertex
\[
V_{2,2d_2}=(0,2(2d-d_2+1)).
\]
Note that the segment $\overline{V_{2,d_1}V_{2,2d_2}}$ lies on the straight line 
$L_2:2(d_1-d_2)x+d_1y=2d_1(2d+1-d_2)$, while the segment $\overline{V_{2,2d}V_{2,2d_2}}$ 
lies on the straight line $L_3:(2d-d_2)x+dy=2d(2d-d_2+1)$. 
Moreover, the slope of $L_3$ is larger than that of $L_1$ owing to 
$-(2d-d_2)/d=-2+d_2/d>-2$.
For $2d_2<l<d_1$, our analysis from Case 1 shows that the vertices of the Newton 
diagram of $\mathcal{F}_{4d+1-l}$ are either on the line $L_2$ or to the right of the 
line $L_2$. Thus, we obtain the configuration of the support for the vector field 
$b(\mathcal{X})$, as shown in Figure \ref{Fig3}(a).
And the Newton diagram of $b(\mathcal{X})$ includes two exterior vertices $V_{1,2d}$ 
and $V_{2,2d_2}$, an inner vertex $V_{2,2d}$, and two edges of type $(2d-d_2, d)$ 
and $(1,1)$, as illustrated in Figure \ref{Fig3}(b).

\begin{figure}[htp]
\centering
(a) \hfil (b)
\caption{(a) Configuration of the support of $b(\mathcal{X})$,
(b)  Newton diagram for $b(\mathcal{X})$.}
\label{Fig3} %\label{Fig3.sub.1} \label{Fig3.sub.2}
\end{figure}

Based on Case 1, we can directly conclude that the Hamiltonian associated with the 
edge of type $(1,1)$ is
\[
h_{2(d+1)}^{(1,1)}=\frac{A_{2d_{\mathbf{s}}}du^{2d}}{2(d+1)}(u^2+v^2).
\]
And the Hamiltonian associated with the edge of type $(2d-d_2, d)$ is
\[
h_{2d(2d-d_2+1)}^{(2d-d_2,d)}
=\frac{d_2v^2}{2(2d-d_2+1)}\left(A_{2d_{\mathbf{s}}}u^{2d}+Bv^{4d-2d_2}\right).
\]
It follows that the Hamiltonians $h_{2d+2}^{(1,1)}$ and $h_{2d(2d-d_2+1)}^{(2d-d_2,d)}$ 
have no factor of the form $v^{t_1}-\lambda u^{t_2}$ with $\lambda\neq 0$.

For $s_1>s_2$, similar results can be obtained only by interchanging $x$ and $y$.
This completes the proof.
\end{proof}

\begin{thebibliography}{99}
\bibitem[1]{bib 1} A. Algaba, E. Gamero, C. Garc\'ia;
The integrability problem for a class of planar systems,
\emph{Nonlinearity}, \textbf{22} (2009), 395-420.
\bibitem[2]{bib 2} A. Algaba, C. Garc\'ia, M. Reyes;
Characterization of a monodromic singular point of a planar vector field,
\emph{Nonlinear Anal.}, \textbf{74} (2011), 5402-5414.
\bibitem[3]{bib 3} A. A. Andronov, E. A. Leontovich, I. I. Gordon, A. G. Ma\u{i}er;
Qualitative Theory of Second-Order Dynamic Systems, Halsted Press, 1973.
\bibitem[4]{bib 4} H. Bass, E. H. Connell, D. Wright;
The Jacobian conjecture: reduction of degree and formal expansion of the inverse,
\emph{Bull. Amer. Math. Soc.}, \textbf{7} (1982), 287-330.
\bibitem[5]{bib 28} F. Braun, F. Fernandes, I. S. Meza-Sarmiento;
Foliations raised by quadratic-like polynomial submersions on the real plane and a sharp result on the real Jacobian conjecture,
\emph{J. Differ. Equ.}, \textbf{449} (2025), 113742, 27pp.
\bibitem[6]{bib 5} F. S. Berezovskaya, N. B. Medvedeva;
A complicated singular point of “center-focus” type and the Newton diagram,
\emph{Sel. Math.}, \textbf{13} (1994), 1-15.
\bibitem[7]{bib 6} F. Braun, J. R. dos Santos Filho;
The real Jacobian conjecture on $\mathbb{R}^2$ is true when one of the components has degree 3,
\emph{Discrete Contin. Dyn. Syst.}, \textbf{26} (2010), 75-87.
\bibitem[8]{bib 7} F. Braun, J. Gin\'{e}, J. Llibre;
A sufficient condition in order that the real Jacobian conjecture in $\mathbb{R}^2$ holds,
\emph{J. Differ. Equ.}, \textbf{260} (2016), 5250-5258.
\bibitem[9]{bib 8} F. Braun, J. Llibre;
A new qualitative proof of a result on the real Jacobian conjecture,
\emph{An. Acad. Bras. Ci\^{e}nc.}, \textbf{87} (2015), 1519-1524.
\bibitem[10]{bib 9} F. Braun, B. Or\'{e}fice-Okamoto;
On polynomial submersions of degree 4 and the real Jacobian conjecture in $\mathbb{R}^2$,
\emph{J. Math. Anal. Appl.}, \textbf{443} (2016), 688-706.
\bibitem[11]{bib 10} F. Braun, C. Valls;
A weight-homogenous condition to the real Jacobian conjecture in $\mathbb{R}^2$,
\emph{Proc. Edinb. Math. Soc.}, (2), \textbf{64} (2021), 1028-1036.
\bibitem[12]{bib 11} A. Cima, A. Gasull, J. Llibre, F. Ma$\Tilde{\rm n}$osas;
Global injectivity of polynomial maps via vector fields, Springer, 1995.
\bibitem[13]{bib 12} A. Cima, A. Gasull, F. Ma$\Tilde{\rm n}$osas;
Injectivity of polynomial local homeomorphisms of $\mathbb{R}^n$,
\emph{Nonlinear Anal.}, \textbf{26} (1996), 877-885.
\bibitem[14]{bib 13} M. Cobo, C. Gutierrez, J. Llibre;
On the injectivity of $C^1$–maps of the real plane,
\emph{Canad. J. Math.}, \textbf{54} (2002), 1187-1201.
\bibitem[15]{bib 14} F. Dumortier, J. Llibre, J. Art\'{e}s;
Qualitative Theory of Planar Differential Systems, Springer, 2006.
\bibitem[16]{bib 15} A. Fernandes, C. Gutierrez, R. Rabanal;
Global asymptotic stability for differentiable vector fields of $\mathbb{R}^2$,
\emph{J. Differ. Equ.}, \textbf{206} (2004), 470-482.
\bibitem[17]{bib 16} L. Gavrilov;
Isochronicity of plane polynomial Hamiltonian systems,
\emph{Nonlinearity}, \textbf{10} (1997), 433-448.
\bibitem[18]{bib 17} J. Gwo\'{z}dziewicz;
The real Jacobian conjecture for polynomials of degree 3,
\emph{Ann. Pol. Math.}, \textbf{76} (2001), 121-125.
\bibitem[19]{bib 18} O.-H. Keller;
Ganze Cremona-Transformationen,
\emph{Monatsh. Math. Phys.}, \textbf{47}(1) (1939), 299-306.
\bibitem[20]{bib 19} C. Liu, Y. Tian;
Sufficient conditions for the $n$-dimensional real Jacobian conjecture,
\emph{J. Differ. Equ.}, \textbf{446} (2025), 113750.
\bibitem[21]{bib 20} J. Llibre, C. Valls;
A sufficient condition for the real Jacobian conjecture in $\mathbb{R}^2$,
\emph{Nonlinear Anal. Real World Appl.}, \textbf{60} (2021), 103298.
\bibitem[22]{bib 21} S. Pinchuk;
A counterexample to the strong real Jacobian conjecture,
\emph{Math. Z.}, \textbf{217} (1994), 1-4.
\bibitem[23]{bib 22} J. D. Randall;
The real Jacobian problem,
\emph{Amer. Math. Soc.}, 1983.
\bibitem[24]{bib 23} M. Sabatini;
A connection between isochronous Hamiltonian centres and the Jacobian conjecture,
\emph{Nonlinear Anal.}, \textbf{34} (1998), 829-838.
\bibitem[25]{bib 24} S. Smale;
Mathematical problems for the next century,
\emph{Math. Intelligencer.}, \textbf{20} (1998), 7-15.
\bibitem[26]{bib 25} Y. Tian, X. Cen;
New sufficient condition for the two-dimensional real Jacobian conjecture 
through the Newton diagram,
\emph{J. Differ. Equ.}, \textbf{410} (2024), 171-196.
\bibitem[27]{bib 26} A. van den Essen;
Polynomial automorphisms and the Jacobian conjecture,
\emph{Birkh\"{a}user Verlag}, Basel, 2000.
\bibitem[28]{bib 27} A. van den Essen, S. Kuroda, A. J. Crachiola;
Polynomial automorphisms and the Jacobian conjecture-new results from the beginning 
of the 21st century, Springer, 2021.
\end{thebibliography}
\end{document}
